\documentclass[12pt,a4paper]{article}

% ── Packages ──────────────────────────────────────────────────────────────────
\usepackage[margin=1in]{geometry}
\usepackage{graphicx}
\usepackage{amsmath, amssymb}
\usepackage{booktabs}
\usepackage{xcolor}
\usepackage{listings}
\usepackage{hyperref}
\usepackage{tikz}
\usepackage{fontenc}
\usepackage{inputenc}
\usepackage{titlesec}
\usepackage{fancyhdr}
\usepackage{mdframed}
\usepackage{enumitem}
\usepackage{multicol}
\usepackage{caption}
\usepackage{float}
\usepackage{array}
\usepackage{tabularx}
\usepackage[most]{tcolorbox}

% ── Colors ────────────────────────────────────────────────────────────────────
\definecolor{teslaRed}{RGB}{204, 0, 0}
\definecolor{codeGray}{RGB}{245, 245, 245}
\definecolor{codeBorder}{RGB}{200, 200, 200}
\definecolor{darkBlue}{RGB}{20, 40, 100}
\definecolor{accentBlue}{RGB}{0, 102, 204}

% ── Hyperref Setup ────────────────────────────────────────────────────────────
\hypersetup{
    colorlinks=true,
    linkcolor=darkBlue,
    urlcolor=accentBlue,
    citecolor=darkBlue,
    pdftitle={Multimodal Self-Supervised Vehicle Log Model},
    pdfauthor={Aviral Srivastava}
}

% ── Code Listing Style ────────────────────────────────────────────────────────
\lstdefinestyle{pythonstyle}{
    language=Python,
    backgroundcolor=\color{codeGray},
    basicstyle=\ttfamily\small,
    keywordstyle=\color{accentBlue}\bfseries,
    stringstyle=\color{teslaRed},
    commentstyle=\color{gray}\itshape,
    numberstyle=\tiny\color{gray},
    numbers=left,
    numbersep=8pt,
    frame=single,
    rulecolor=\color{codeBorder},
    breaklines=true,
    showstringspaces=false,
    tabsize=4
}
\lstdefinestyle{bashstyle}{
    language=bash,
    backgroundcolor=\color{black!90},
    basicstyle=\ttfamily\small\color{white},
    breaklines=true,
    frame=none,
    showstringspaces=false
}

% ── Header / Footer ───────────────────────────────────────────────────────────
\pagestyle{fancy}
\fancyhf{}
\fancyhead[L]{\small\textcolor{teslaRed}{\textbf{Tesla Assessment}}}
\fancyhead[R]{\small\textcolor{gray}{Aviral Srivastava}}
\fancyfoot[C]{\thepage}
\renewcommand{\headrulewidth}{0.5pt}

% ── Section Formatting ────────────────────────────────────────────────────────
\titleformat{\section}{\large\bfseries\color{darkBlue}}{}{0em}{\thesection.\quad}[\color{teslaRed}\titlerule]
\titleformat{\subsection}{\normalsize\bfseries\color{darkBlue}}{}{0em}{\thesubsection.\quad}

% ══════════════════════════════════════════════════════════════════════════════
\begin{document}

% ── Title Page ────────────────────────────────────────────────────────────────
\begin{titlepage}
    \centering
    \vspace*{1cm}

    {\Huge\bfseries\color{darkBlue} Multimodal Self-Supervised\\[0.3em]
    \textcolor{teslaRed}{Vehicle Log Model}}\\[1.5cm]

    {\large\textbf{Technical Documentation \& Architecture Overview}}\\[0.5cm]
    \rule{0.6\textwidth}{0.5pt}\\[1cm]

    {\large\textbf{Author:} Aviral Srivastava}\\[0.3cm]
    {\large\textbf{Date:} \today}\\[2cm]

    \begin{tcolorbox}[colback=darkBlue!5, colframe=darkBlue, width=0.85\textwidth,
        title={\textbf{Abstract}}, fonttitle=\bfseries\color{white},
        coltitle=white, colbacktitle=darkBlue]
        This document presents a self-supervised multimodal foundation model for
        autonomous vehicle log understanding. The model fuses five heterogeneous
        sensor modalities — camera (BEV), LiDAR, CAN bus, GPS, and HD-Map —
        using a DINOv2-style contrastive pretraining strategy with an EMA teacher
        network. The architecture is designed for efficient transfer to downstream
        tasks such as anomaly detection, driver behaviour classification, and
        trajectory prediction.
    \end{tcolorbox}

    \vfill
    \small\textcolor{gray}{Tesla Assessment | Self-Supervised Learning | PyTorch}
\end{titlepage}

\tableofcontents
\newpage

% ══════════════════════════════════════════════════════════════════════════════
\section{Introduction}

Modern autonomous driving systems generate enormous volumes of heterogeneous
sensor data — LiDAR point clouds, camera images, GPS traces, CAN bus signals, and
HD-Map graphs. Supervised learning on such data is expensive due to annotation costs,
while fully supervised models generalise poorly across driving environments.

This project proposes a \textbf{self-supervised pretraining framework} that learns
rich, transferable representations from unlabelled vehicle logs. The key contributions are:

\begin{itemize}[leftmargin=1.5em]
    \item A unified multimodal pipeline ingesting 5 sensor modalities from HDF5 logs.
    \item A contrastive SSL objective combining NT-Xent and DINO-style losses.
    \item A Fusion Transformer that produces a single 2048-dimensional embedding.
    \item Efficient data augmentation including BEV multi-crop, LiDAR dropout, and CAN/GPS jitter.
    \item Downstream evaluation via linear probes and trajectory heads.
\end{itemize}

% ══════════════════════════════════════════════════════════════════════════════
\section{System Architecture}

\subsection{Overview}

\begin{figure}[H]
\centering
\begin{tcolorbox}[colback=codeGray, colframe=codeBorder, width=\textwidth]
\begin{verbatim}
  Raw Logs (camera / LiDAR / CAN / GPS / HD-Map)
                      |
                      v
     ┌────────────────────────────────────┐
     │  VehicleLogDataset  (pipeline.py)  │
     │  BEV projection  |  LiDAR voxels  │
     │  CAN sequences   |  GPS tensors    │
     │  PyG lane graph                    │
     └──────────────┬─────────────────────┘
                    │  DualViewAugment
            ┌───────┴───────┐
          View 1          View 2
            │               │
  ┌─────────▼───────────────▼─────────┐
  │    Parallel Modality Encoders      │
  │  BEV   → DINOv2 ViT-B/14  → 768  │
  │  LiDAR → PointNet         → 1024  │
  │  CAN   → Bi-LSTM          →  512  │
  │  GPS   → Transformer      →  512  │
  │  Map   → GCNConv (3-lay)  →  512  │
  └────────────────┬──────────────────┘
                   │  ProjectionHead (MLP, BN, GELU)
                   v  1024-dim L2-normed
  ┌────────────────────────────────────┐
  │  FusionTransformer                 │
  │  5 tokens × 1024 → cross-attn     │
  │  → flatten → Linear → 2048-dim    │
  └───────────────┬────────────────────┘
                  │
  ┌───────────────▼────────────────────┐
  │  Contrastive Losses                │
  │  NT-Xent (per modality + fused)    │
  │  DINO-style vs EMA teacher         │
  └────────────────────────────────────┘
\end{verbatim}
\end{tcolorbox}
\caption{End-to-end architecture of the multimodal SSL pipeline.}
\end{figure}

\subsection{Modality Encoders}

Each sensor modality is processed by a dedicated encoder:

\begin{table}[H]
\centering
\renewcommand{\arraystretch}{1.4}
\begin{tabularx}{\textwidth}{lXll}
\toprule
\textbf{Modality} & \textbf{Encoder} & \textbf{Output Dim} & \textbf{Input Shape} \\
\midrule
Camera / BEV & DINOv2 ViT-B/14 & 768 & $(3, 224, 224)$ \\
LiDAR         & PointNet MLP     & 1024 & $(N, 3)$ pts \\
CAN bus       & Bi-LSTM (2-layer) & 512 & $(T, D)$ signals \\
GPS           & Transformer (4-head) & 512 & $(T, 3)$ coords \\
HD-Map        & GCNConv (3-layer) & 512 & Graph $(M, 8)$ \\
\bottomrule
\end{tabularx}
\caption{Summary of per-modality encoders and output dimensionalities.}
\end{table}

\subsection{Fusion Transformer}

After encoding, each modality embedding is projected to a unified 1024-dim space via a
\texttt{ProjectionHead} (3-layer MLP with BatchNorm and GELU activations, L2-normalised).
The five projected tokens are then fused using a \textbf{FusionTransformer}:

\begin{equation}
    \mathbf{z} = \text{Linear}\bigl(\text{Flatten}(\text{TransformerEncoder}([\mathbf{p}_1, \mathbf{p}_2, \mathbf{p}_3, \mathbf{p}_4, \mathbf{p}_5]))\bigr) \in \mathbb{R}^{2048}
\end{equation}

where $\mathbf{p}_i$ are the projected modality tokens.

\subsection{EMA Teacher Network}

Following DINO, an \textbf{Exponential Moving Average (EMA) teacher} is maintained:

\begin{equation}
    \theta_t \leftarrow m \cdot \theta_t + (1-m) \cdot \theta_s, \quad m = 0.996
\end{equation}

The teacher provides stable target embeddings without backpropagation, preventing
representation collapse during contrastive learning.

% ══════════════════════════════════════════════════════════════════════════════
\section{Data Pipeline (\texttt{pipeline.py})}

\subsection{Dataset Format}

Input data is stored in \textbf{HDF5} format with one group per driving segment:

\begin{lstlisting}[style=pythonstyle, caption={HDF5 dataset structure}]
logs.h5
├── segment_0000/
│   ├── lidar       (N_pts, 4)   float32  # x, y, z, intensity
│   ├── camera      (3, H, W)    float32  # RGB image
│   ├── can         (T, 50)      float32  # CAN bus signals
│   ├── gps         (T, 3)       float32  # lat, lon, velocity
│   ├── map_nodes   (M, 8)       float32  # lane graph nodes
│   └── map_edges   (2, E)       int64    # lane graph edges
└── segment_0001/ ...
\end{lstlisting}

\subsection{VehicleLogDataset}

The \texttt{VehicleLogDataset} class (\texttt{pipeline.py}) handles:
\begin{itemize}[leftmargin=1.5em]
    \item Loading HDF5 segments lazily (memory-efficient).
    \item BEV projection of LiDAR via \texttt{project\_to\_bev()} — 200$\times$200 grid, 50m range.
    \item Padding/truncating sequences to fixed length (\texttt{seq\_len=50}).
    \item Converting map data to PyTorch Geometric \texttt{Data} objects.
    \item Applying \texttt{DualViewAugment} to produce two augmented views per sample.
\end{itemize}

\subsection{Data Augmentation Strategy}

\begin{table}[H]
\centering
\renewcommand{\arraystretch}{1.3}
\begin{tabular}{lll}
\toprule
\textbf{Modality} & \textbf{Augmentation} & \textbf{Parameters} \\
\midrule
BEV (Global)  & RandomResizedCrop, HFlip, ColorJitter & scale $(0.4, 1.0)$ \\
BEV (Local)   & RandomResizedCrop, GaussianBlur       & scale $(0.05, 0.4)$, 6 crops \\
LiDAR         & Point dropout + Gaussian jitter        & 30\% drop, $\sigma=0.02$ \\
CAN bus       & Gaussian noise + temporal masking      & $\sigma=0.01$, 10\% mask \\
GPS           & Coordinate + velocity jitter           & $\sigma_{coord}=0.0001$ \\
\bottomrule
\end{tabular}
\caption{Multi-modal augmentation pipeline.}
\end{table}

\subsection{Custom Collate Function}

The \texttt{ssl\_collate\_fn} stacks tensors per modality across batch items and
batches PyG graphs using \texttt{Batch.from\_data\_list()} for efficient GPU parallelism.

% ══════════════════════════════════════════════════════════════════════════════
\section{Training Procedure (\texttt{train.py})}

\subsection{Loss Functions}

The total SSL loss combines three objectives:

\begin{equation}
    \mathcal{L}_{\text{total}} = \mathcal{L}_{\text{NT-Xent}}^{\text{fused}} + \sum_{m} \mathcal{L}_{\text{NT-Xent}}^{(m)} + \mathcal{L}_{\text{DINO}}
\end{equation}

\textbf{NT-Xent (Normalised Temperature-scaled Cross Entropy):}
\begin{equation}
    \mathcal{L}_{\text{NT-Xent}} = -\log \frac{\exp(\mathbf{z}_i \cdot \mathbf{z}_j / \tau)}{\sum_{k \neq i} \exp(\mathbf{z}_i \cdot \mathbf{z}_k / \tau)}, \quad \tau = 0.1
\end{equation}

\textbf{DINO Loss} uses Sinkhorn-Knopp normalisation on teacher soft labels:
\begin{equation}
    \mathcal{L}_{\text{DINO}} = -\sum_k p_t^{(k)} \log p_s^{(k)}
\end{equation}

\subsection{Optimiser and Learning Rate Schedule}

\begin{itemize}[leftmargin=1.5em]
    \item \textbf{Optimiser:} AdamW with decoupled weight decay ($\beta_1=0.9, \beta_2=0.999$)
    \item \textbf{LR Schedule:} Cosine annealing with 5\% linear warmup
    \item \textbf{Gradient clipping:} $\|\nabla\|_2 \leq 3.0$
    \item \textbf{Mixed Precision:} AMP via \texttt{torch.amp.GradScaler}
\end{itemize}

\begin{equation}
    \eta(t) = \eta_{\min} + \frac{1}{2}(\eta_{\max} - \eta_{\min})\left(1 + \cos\left(\frac{\pi (t - t_w)}{T - t_w}\right)\right)
\end{equation}

\subsection{Key Hyperparameters}

\begin{table}[H]
\centering
\renewcommand{\arraystretch}{1.3}
\begin{tabular}{lll}
\toprule
\textbf{Parameter} & \textbf{Default} & \textbf{Description} \\
\midrule
\texttt{lr}             & $1 \times 10^{-4}$  & AdamW base learning rate \\
\texttt{batch\_size}    & 8                    & Per-GPU batch size \\
\texttt{temperature}    & 0.1                  & NT-Xent temperature $\tau$ \\
\texttt{ema\_momentum}  & 0.996                & Teacher EMA coefficient \\
\texttt{epochs}         & 20                   & Training epochs \\
\texttt{weight\_decay}  & 0.04                 & AdamW weight decay \\
\texttt{freeze\_dino}   & True                 & Freeze DINOv2 backbone \\
\texttt{amp}            & False                & Mixed precision training \\
\bottomrule
\end{tabular}
\caption{Training hyperparameters.}
\end{table}

\subsection{Checkpoint Strategy}

Checkpoints are saved after every epoch as \texttt{ckpt\_epoch\{N\}.pt} and a rolling
\texttt{ckpt\_last.pt}. Each checkpoint stores student weights, teacher weights,
optimiser state, and scaler state for seamless resumption via \texttt{--resume}.

% ══════════════════════════════════════════════════════════════════════════════
\section{Evaluation (\texttt{eval.py})}

\subsection{Linear Probe}

A linear classifier is trained on frozen representations for binary classification tasks
(e.g., accident detection, driver behaviour):

\begin{equation}
    \hat{y} = \text{softmax}(W \mathbf{z} + b)
\end{equation}

\textbf{Target metric:} AUROC $> 0.90$ for accident prediction.

\subsection{Trajectory Prediction}

A lightweight MLP head predicts 30-step future trajectories:

\begin{itemize}[leftmargin=1.5em]
    \item \textbf{ADE} (Average Displacement Error): $< 1.5$ m
    \item \textbf{FDE} (Final Displacement Error): $< 3.0$ m
\end{itemize}

\subsection{Expected Results}

\begin{table}[H]
\centering
\renewcommand{\arraystretch}{1.3}
\begin{tabular}{lll}
\toprule
\textbf{Task} & \textbf{Metric} & \textbf{Target} \\
\midrule
Accident prediction   & AUROC & $>0.90$ \\
Driver behaviour      & AUROC & $>0.85$ \\
Trajectory (30-step)  & ADE   & $<1.5$ m \\
Trajectory (30-step)  & FDE   & $<3.0$ m \\
\bottomrule
\end{tabular}
\caption{Downstream task evaluation targets.}
\end{table}

% ══════════════════════════════════════════════════════════════════════════════
\section{File Structure}

\begin{lstlisting}[style=pythonstyle, caption={Project directory layout}]
tesla-assessment/
├── train.py          # Main SSL training entry point
├── eval.py           # Linear probe + trajectory evaluation
├── pipeline.py       # Dataset, DataLoader, collate_fn
├── requirements.txt  # Python dependencies
├── README.md         # Project overview
├── models/
│   ├── __init__.py
│   ├── encoders.py   # 5 modality encoders + ProjectionHead
│   ├── fusion.py     # FusionTransformer + EMA utils
│   └── losses.py     # NT-Xent, Sinkhorn-Knopp, DINO loss
└── utils/
    ├── __init__.py
    ├── bev.py        # BEV projection + LiDAR voxelization
    └── augment.py    # Multi-crop, dropout, CAN/GPS jitter
\end{lstlisting}

% ══════════════════════════════════════════════════════════════════════════════
\section{Quick Start}

\subsection{Installation}

\begin{lstlisting}[style=bashstyle, caption={Environment setup}]
conda create -n vssl python=3.10
conda activate vssl
pip install -r requirements.txt

# PyG (match CUDA version)
pip install torch-geometric torch-scatter torch-sparse \
    -f https://data.pyg.org/whl/torch-2.2.0+cu121.html
\end{lstlisting}

\subsection{Training}

\begin{lstlisting}[style=bashstyle, caption={Launch SSL pretraining}]
python train.py \
    --data_root /path/to/logs_256.h5 \
    --epochs 20 \
    --batch_size 8 \
    --num_workers 0 \
    --freeze_dino \
    --log_every 4
\end{lstlisting}

\subsection{Google Colab (Recommended)}

For faster training on cloud GPU (T4), use the provided \texttt{Tesla\_Colab\_Training.ipynb}
notebook which automatically generates a synthetic dataset matching the HDF5 format and
trains the model end-to-end in approximately 15 minutes.

% ══════════════════════════════════════════════════════════════════════════════
\section{Dependencies}

\begin{table}[H]
\centering
\renewcommand{\arraystretch}{1.3}
\begin{tabular}{ll}
\toprule
\textbf{Package} & \textbf{Purpose} \\
\midrule
\texttt{torch $\geq$ 2.0}       & Deep learning framework \\
\texttt{torchvision}             & Image transforms \\
\texttt{torch-geometric}         & Graph neural networks (HD-Map) \\
\texttt{h5py}                    & HDF5 dataset loading \\
\texttt{numpy, pandas}           & Data processing \\
\texttt{tensorboard}             & Training visualisation \\
\texttt{wandb}                   & Experiment tracking (optional) \\
\bottomrule
\end{tabular}
\caption{Key project dependencies.}
\end{table}

% ══════════════════════════════════════════════════════════════════════════════
\section{Conclusion}

This project demonstrates a complete self-supervised multimodal learning pipeline for
autonomous vehicle data. By combining DINOv2-style contrastive pretraining with a
Fusion Transformer, the system learns semantically rich representations from unlabelled
sensor logs that transfer effectively to safety-critical downstream tasks.

The modular design — with separate encoder, fusion, loss, and pipeline components —
enables straightforward extension to additional modalities or alternative SSL objectives
such as MAE or VICReg.

\vspace{1em}
\begin{tcolorbox}[colback=darkBlue!5, colframe=darkBlue]
\textbf{GitHub Repository:}
\href{https://github.com/Aviralsriv/tesla-assesment-2}{\texttt{github.com/Aviralsriv/tesla-assesment-2}}
\end{tcolorbox}

% ══════════════════════════════════════════════════════════════════════════════
\end{document}
